% Reproducibility appendix. Requires amsmath; citation key: OpenAI4.
\section{Exact coefficient certificates and reproducibility}
\label{app:exact-certificates}

This appendix specifies the finite arithmetic certificates used in the proof.
The three-state inequalities are verified by two different polynomial
constructions.  The four-state inequalities use the published certificate
of \cite{OpenAI4}, whose two complete traversals were rerun.
All sign decisions use integer or rational arithmetic.

\subsection{The three-state residuals and their complete domain}

For \(p\in\Delta_3\), retain the definitions
\[
 x(p)=\frac16\sum_i(3p_i-1)^2,\qquad
 y(p)=\frac16\sum_i(3p_i-1)^3,\qquad
 H=y-\frac{x^2}{2},\quad F=\frac32x-y,\quad G=1-x.
\]
For \(\pi\in S_3\), put
\[
 z_\pi=3\sum_i p_iq_{\pi(i)},\qquad
 r_{\pi,i}=\frac{3p_iq_{\pi(i)}}{z_\pi},\qquad
 \epsilon=\frac1{1000}.
\]
A term \(z_\pi f(r_\pi)\) is defined as zero when \(z_\pi=0\).
The two residuals certified below are
\begin{align}
 R_A(p,q)
 &=x(p)+x(q)-\frac16\sum_{\pi\in S_3}z_\pi x(r_\pi)\notag\\
 &\quad-H(p)F(q)-F(p)H(q)-\epsilon x(p)x(q), \label{eq:app-RA}\\
 R_B(p,q)
 &=\frac16\sum_{\pi\in S_3}z_\pi H(r_\pi)\notag\\
 &\quad-H(p)G(q)-G(p)H(q)-\epsilon x(p)x(q). \label{eq:app-RB}
\end{align}

It suffices to work in the ordered chambers for \(p\) and \(q\):
all scalar functions are symmetric, and the fusion terms average over
every relative permutation.  Introduce nonnegative variables
\((a,b,c)\), \((d,e,f)\), and let \(S=a+b+c>0\), \(T=d+e+f>0\).  Set
\[
 p=\frac{(6a+3b+2c,\;3b+2c,\;2c)}{6S},
 \qquad
 q=\frac{(6d+3e+2f,\;3e+2f,\;2f)}{6T}.
\]
This covers the whole ordered simplex.  Indeed, for
\(p_1\ge p_2\ge p_3\), the normalized coordinates
\[
 (a,b,c)=(p_1-p_2,\;2(p_2-p_3),\;3p_3)
\]
are nonnegative and sum to one.  Thus the certificate is not restricted
to sampled posterior vectors or to a neighborhood of the uniform law.

\subsection{An explicit integer polynomial certificate}

Define
\[
 X=4a^2+2ab+b^2,\qquad
 Y=8a^3+6a^2b-3ab^2-b^3,
\]
and use primes for the same expressions in \((d,e,f)\).  Direct
substitution gives
\[
 x(p)=\frac{X}{4S^2},\qquad y(p)=\frac{Y}{8S^3}.
\]
Writing
\[
 \widehat H=4YS-X^2,\qquad
 \widehat F=3XS-Y,\qquad
 \widehat G=4S^2-X,
\]
we have
\[
 H(p)=\frac{\widehat H}{32S^4},\qquad
 F(p)=\frac{\widehat F}{8S^3},\qquad
 G(p)=\frac{\widehat G}{4S^2}.
\]
The auxiliary signs used in the expectation argument follow, for example,
from the explicit expansions
\begin{align*}
 \widehat F
 &=4a^3+12a^2b+12ab^2+4b^3+12a^2c+6abc+3b^2c,\\
 \widehat G
 &=6ab+8ac+3b^2+8bc+4c^2.
\end{align*}

Let \(P=(6a+3b+2c,3b+2c,2c)\) and
\(Q=(6d+3e+2f,3e+2f,2f)\).  For each permutation define
\[
 U_{\pi,i}=P_iQ_{\pi(i)},\quad
 Z_\pi=\sum_iU_{\pi,i},\quad
 W_{\pi,i}=3U_{\pi,i}-Z_\pi,
\]
\[
 V_\pi=\sum_iW_{\pi,i}^{\,2},\qquad
 J_\pi=36Z_\pi W_{\pi,1}W_{\pi,2}W_{\pi,3}-V_\pi^2.
\]
Since \(\sum_iW_{\pi,i}=0\), the weighted fusion terms satisfy
\[
 z_\pi=\frac{Z_\pi}{12ST},\qquad
 x(r_\pi)=\frac{V_\pi}{6Z_\pi^2},\qquad
 H(r_\pi)=\frac{J_\pi}{72Z_\pi^4}.
\]
In particular, the factor \(z_\pi\) from the posterior product law reduces
the denominator powers to one and three.

Set \(D_A=\prod_\pi Z_\pi\), \(D_B=\prod_\pi Z_\pi^3\), and define
\begin{align*}
 K_A={}&8000\bigl(XS^2T^4+X'S^4T^2\bigr)
       -125\bigl(\widehat H\widehat F'T+\widehat F\widehat H'S\bigr)
       -2XX'S^2T^2,\\
 K_B={}&-125\bigl(\widehat H\widehat G'T^2+
                         \widehat G\widehat H'S^2\bigr)-XX'S^2T^2.
\end{align*}
The following two polynomials have integer coefficients:
\begin{align}
 \mathcal P_A
 &=27K_AD_A-2000S^3T^3
       \sum_\pi V_\pi\prod_{\rho\ne\pi}Z_\rho, \label{eq:app-PA}\\
 \mathcal P_B
 &=81K_BD_B+250S^3T^3
       \sum_\pi J_\pi\prod_{\rho\ne\pi}Z_\rho^3. \label{eq:app-PB}
\end{align}
Their exact relation to the desired residuals is
\begin{equation}
 \mathcal P_A=864000S^4T^4D_A R_A,\qquad
 \mathcal P_B=1296000S^4T^4D_B R_B.                 \label{eq:app-clearing}
\end{equation}

Expanding \eqref{eq:app-PA}--\eqref{eq:app-PB} in the six nonnegative
barycentric variables gives the complete coefficient inventory:
\begin{center}
\begin{tabular}{lrrrrr}
\hline
Certificate & Bidegree & Positions & Positive & Zero & Negative\\
\hline
A & \((10,10)\) & \(4\,356\) & \(3\,873\) & \(483\) & \(0\)\\
B & \((22,22)\) & \(76\,176\) & \(71\,627\) & \(4\,549\) & \(0\)\\
\hline
\end{tabular}
\end{center}
The number of positions is \(\binom{n+2}{2}^2\) in bidegree \((n,n)\).
Every coefficient is checked; no subdivision is needed for these
three-state polynomials.

Nonnegative coefficients imply \(\mathcal P_A,\mathcal P_B\ge0\).
For interior probability vectors, \(Z_\pi>0\) for all \(\pi\), so
\eqref{eq:app-clearing} proves \(R_A,R_B\ge0\).
Boundary points are included by continuity: \(x\) and \(H\) are bounded on
the probability simplex, hence \(z_\pi x(r_\pi)\) and
\(z_\pi H(r_\pi)\) tend to zero when \(z_\pi\) tends to zero.
Approximation by interior pairs therefore establishes
\eqref{eq:app-RA}--\eqref{eq:app-RB} on the entire domain.

\subsection{Independent verification}

The source \path{analytic/q3_exact_pair_simple.cpp} constructs rational
homogeneous polynomials directly from the definitions of the posterior
invariants.  Separately,
\path{computation/q3_independent_coefficients.cpp} expands
\eqref{eq:app-PA}--\eqref{eq:app-PB} in a sparse six-variable integer
representation.  Its construction does not use the first encoder's
invariants, homogenization, tensor multiplication, or rational-polynomial
operations.  Only after constructing its polynomials does it call the
first encoder to compare the complete coefficient arrays.

If \(\mathcal C_A,\mathcal C_B\) denote the rational homogeneous outputs of
the first encoder, every coefficient satisfies the exact identities
\[
 \mathcal P_A=144000\cdot12^6\,\mathcal C_A,\qquad
 \mathcal P_B=216000\cdot12^{18}\,\mathcal C_B.
\]
The comparison uses integer cross-products with the stored denominators.
Both constructions use arbitrary-precision integers.  The independent
implementation packs six exponents into five-bit fields and checks that
each group's total degree is below \(32\); the maximum used is \(22\),
so exponent addition cannot carry between fields.

Both implementations completed with assertions enabled and
UndefinedBehaviorSanitizer.  The clean certificates accept at their root,
so floating-point edge-selection coordinates in the first encoder are
never consulted.  No sampled inequality or floating-point sign decision
is used.

The supplementary files contain both implementations, the complete coefficient tables, and the verification output. The independent program regenerates every coefficient, including all zeros.

\subsection{Four-state certificate}

The four-state pair inequalities and their certificate algorithm are those of \cite{OpenAI4}; the supplementary files include the unchanged verifier from the version linked in that reference.

The verifier works on products of ordered posterior simplices and sums
over all \(24\) relative colour permutations.  It certifies subdomains
using either a Taylor lower bound for reciprocal powers with a
nonnegative multiplier, or cubic interpolation with a rigorously bounded
cleared polynomial remainder.  Exact midpoint bisection covers a
subdomain whenever neither certificate yet applies.  A successful
traversal certifies every resulting subdomain; see \cite{OpenAI4} for the
inequalities and the justification of these two tests.

Both full traversals were executed with the published C++ source unchanged,
exact integer/rational arithmetic, and assertions enabled:
\begin{center}
\begin{tabular}{lrrrr}
\hline
Mode & Nodes & Max.\ depth & Interpolation leaves & Taylor leaves\\
\hline
A & \(4\,479\) & \(19\) & \(2\,074\) & \(166\)\\
B & \(1\,935\) & \(20\) & \(871\) & \(97\)\\
\hline
\end{tabular}
\end{center}
Both exited with status zero.  Their terminal lines were
\begin{verbatim}
done 4479 19 lp=2074 lt=166
done 1935 20 lp=871 lt=97
\end{verbatim}
Complete output records are included with the supplementary code.

Additional checks verify the four homogeneous scalar polynomials' coefficient minima, all \(24\) symmetry actions, and the exact arithmetic, interpolation, and subdivision operations under UndefinedBehaviorSanitizer.  The scalar coefficient minima are
\(19/50,0,87/1000,0\) for \(H+1/2,F,G,(3-\mathrm{Euler})H\),
respectively; these are coefficient minima, not asserted minima of the
functions.  The lightweight checks are supplemented by the two complete traversals reported above.

\subsection{Compact reproduction instructions}

The accompanying source archive preserves the cited files and their
relative paths.  Requirements are Python~3, a C++17 compiler with the
GCC-compatible headers used by the published verifier, and Boost
multiprecision headers.  Keep assertions enabled.  If Boost is outside
the compiler's default include path, set \texttt{CPLUS\_INCLUDE\_PATH}
to its include directory.  From the archive root:
\begin{verbatim}
cd computation
g++ -std=c++17 -O2 -UNDEBUG -fsanitize=undefined \
  -fno-sanitize-recover=all \
  ../analytic/q3_exact_pair_simple.cpp -o q3_check
./q3_check 0
./q3_check 1
g++ -std=c++17 -O2 -UNDEBUG -fsanitize=undefined \
  -fno-sanitize-recover=all \
  q3_independent_coefficients.cpp -o q3_independent
./q3_independent 0
./q3_independent 1
cd published_q4/verification
python3 checks/verify.py
g++ -std=c++17 -O3 -UNDEBUG \
  checks/certificate.cpp -o artifacts/certificate
./artifacts/certificate
./artifacts/certificate b
\end{verbatim}
Every command must exit with status zero.  The q3 outputs must reproduce
the coefficient counts above, and the q4 executions must print the two
stated completion lines. The probabilistic deductions from these algebraic inequalities are supplied in the main
text.

